% ECONOMICS_PROBLEM: {"id": "G12-249", "title": "Excess volatility and return predictability", "jel_code": "G12", "source_id": "OP-0092"}
% !TeX program = xelatex
\documentclass[11pt,a4paper]{article}
\usepackage[margin=23mm]{geometry}
\usepackage{fontspec}
\setmainfont{Latin Modern Roman}
\usepackage{amsmath,amssymb,mathtools}
\usepackage{xcolor,enumitem,booktabs,longtable,array,xurl,hyperref}
\definecolor{muted}{HTML}{53636C}
\hypersetup{colorlinks=true,urlcolor=blue,linkcolor=blue}
\setlength{\parindent}{0pt}
\setlength{\parskip}{5pt}
\newcommand{\R}{\mathbb R}
\newcommand{\E}{\mathbb E}
\newcommand{\N}{\mathbb N}
\newcommand{\ind}{\mathbf 1}
\newcommand{\Var}{\operatorname{Var}}
\newcommand{\Cov}{\operatorname{Cov}}
\newcommand{\supp}{\operatorname{supp}}
\DeclareMathOperator*{\argmin}{arg\,min}
\DeclareMathOperator*{\argmax}{arg\,max}
\newcommand{\entrymeta}[3]{{\sffamily\small\color{muted}\textbf{\texttt{#1}}\quad|\quad #2\par #3\par}}
\newcommand{\Problemref}[1]{\texttt{#1}}
\newcommand{\JELref}[2]{\texttt{#2} (JEL \texttt{#1})}
\begin{document}
\section*{Excess volatility and return predictability}
\textbf{Problem ID:} \texttt{G12-249}\quad \textbf{JEL:} \texttt{G12}\par
% BEGIN ECONOMICS STATEMENT
\label{problem:OP-0092}
\label{audited:ECON-037}
\label{atlas:prob:F2}

\noindent\textbf{Atlas ID:} F2; \textbf{original number:} 92; \textbf{field:} Finance. \textbf{Classification:} Editorial research formulation (theoretical/empirical); not a certified unsolved theorem.

\subsection*{Definitions and assumptions}
On a filtered probability space, a share has ex-dividend price $P_t\ge0$, subsequent dividends $D_{t+1}\ge0$ and a positive one-period stochastic discount factor $m_{t+1}$. Define $M_{t,t+j}=\prod_{k=1}^{j}m_{t+k}$ and suppose $P_t=\mathbb E_t[m_{t+1}(D_{t+1}+P_{t+1})]$ with finite expectations. If discounted dividends have a finite expected sum and $\lim_{H\to\infty}\mathbb E_t[M_{t,t+H}P_{t+H}]=0$, then the no-bubble present-value formula holds. Model $\theta$ specifies investor information, cash-flow law, discounting and an equilibrium selection with feasible asset trades and market clearing.

\subsection*{Formal research question}
Using price, dividend, expectation-survey and position data, identify the components of valuation variation attributable to cash-flow expectations, discount rates and any admissible terminal-value component. Under the stated transversality condition,
\[
 P_t=\mathbb E_t\!\left[\sum_{j\ge1}M_{t,t+j}D_{t+j}\right].
\]
Characterize observationally equivalent models giving different decompositions and determine what additional information separates them. Test return predictability prospectively using a fixed forecasting information set and loss function. If a bubble component is allowed, specify its law, integrability and equilibrium feasibility rather than labeling unexplained residual variation a bubble.

\subsection*{Interpretation and scope}
With deterministic discount $q\in(0,1)$, let $V_t^*=\sum_{j\ge1}q^jD_{t+j}$ be square integrable. If $P_t=\mathbb E_t[V_t^*]$, the law of total variance gives $\operatorname{Var}(P_t)\le\operatorname{Var}(V_t^*)$ whenever those unconditional variances exist. This bound is not a universal restriction under stochastic discounting, changing cash-flow definitions or nonstationary series. Log-linear present-value formulas are approximations with remainder and terminal terms; estimation should assess these rather than treating them as exact identities. Predictability may reflect changing risk compensation and need not imply feasible arbitrage. Survey measures require a model linking reported beliefs to price-setting investors. The share’s dividends must include the chosen treatment of buybacks, issuance and liquidation. Altering payout definitions changes the valuation object. A valid decomposition also states the investor filtration; aggregate survey averages need not equal the relevant conditional expectation.

\subsection*{Source audit and status}
[S1] supplies the historical variance-bound benchmark. The atlas’s ``Campbell–Shiller (1988)'' is ambiguous: [S2] selects the dividend-price article and [S3] records the alternative joint paper. [S4] reviews discount-rate variation. This entry is an editorial identification agenda, not a claim that established present-value identities remain unproved.

\subsection*{References}

\begin{enumerate}[label={[S\arabic*]},leftmargin=*]

\item Robert J. Shiller (1981). \emph{Do Stock Prices Move Too Much to Be Justified by Subsequent Changes in Dividends?}. American Economic Review 71(3), 421–436. \url{https://www.aeaweb.org/aer/top20/71.3.421-436.pdf}.
\textit{Locator and role:} Publisher or author abstract / bibliographic record; Variance-bound evidence under present-value benchmark assumptions. Provides background or a benchmark result; does not establish that the editorial question remains unresolved in 2026.

\item John Y. Campbell and Robert J. Shiller (1988). \emph{The Dividend-Price Ratio and Expectations of Future Dividends and Discount Factors}. Review of Financial Studies 1(3), 195–228. \url{https://doi.org/10.1093/rfs/1.3.195}.
\textit{Locator and role:} Publisher or author abstract / bibliographic record; Log-linear dividend-price present-value decomposition. Two relevant joint papers appeared in 1988; the title was absent. This explicitly selected candidate is accompanied by the alternative Stock Prices, Earnings, and Expected Dividends.

\item John Y. Campbell and Robert J. Shiller (1988). \emph{Stock Prices, Earnings, and Expected Dividends}. Journal of Finance 43(3), 661–676. \url{https://doi.org/10.1111/j.1540-6261.1988.tb04598.x}.
\textit{Locator and role:} Publisher or author abstract / bibliographic record; Alternative 1988 joint paper on stock-price and dividend expectations. Provides background or a benchmark result; does not establish that the editorial question remains unresolved in 2026.

\item John H. Cochrane (2011). \emph{Presidential Address: Discount Rates}. Journal of Finance 66(4), 1047–1108. \url{https://doi.org/10.1111/j.1540-6261.2011.01671.x}.
\textit{Locator and role:} Publisher or author abstract / bibliographic record; Discount-rate variation and return predictability review. Provides background or a benchmark result; does not establish that the editorial question remains unresolved in 2026.

\end{enumerate}

\subsection*{Integrated refinement: ECON-037}
{\sffamily\small\color{muted}High-dimensional investor learning and return predictability\par}
This refinement adopts the additional source conventions
(page~\pageref{audited:conventions}), subject to its local restrictions.
\textbf{High-dimensional investor learning as a predictability mechanism.}
Let investors forecast from $X_t\in\mathbb R^{p_T}$ through a predeclared learning
rule and coefficients $\widehat b_{it}$; let the objective conditional return mean
be $m_P(X_t)$. Use
\[
D_t=\operatorname{Var}_i(X_t^{\mathsf T}\widehat b_{it}),\qquad
B_t=\mathbb E_i[X_t^{\mathsf T}\widehat b_{it}]-m_P(X_t).
\]
Identify rules jointly accounting for disagreement, bias and subsequent excess
return predictability as $p_T$ grows. Match forecast horizons and information
sets, specify signal scaling and growth, moments and measurement error, and
evaluate on later held-out observations. Distinguish private information and
measurement error from belief disagreement. If learning is a causal asset-pricing
explanation, supply preferences, portfolio constraints, market clearing and the
belief-to-price mechanism; predictive fit alone is insufficient.
\subsection*{Supplied source audit for ECON-037}
Local [S1], [S2], etc. in this audit and the references immediately below
belong to ECON-037, independently of the original entry's reference numbering.
[S1], PDF pp. 17 and 19--21, supports priors and high-dimensional learning. These disagreement/bias diagnostics and rule-selection criteria are editorial.
\subsection*{References for integrated refinement ECON-037}
{\footnotesize\raggedright
\par\noindent\textbf{[S1]} Stefan Nagel (2021). \emph{Subjective Beliefs in Asset Pricing}. Innsbruck lecture slides, June 2021; PDF has two slides on many pages.\par \textit{Verified locator / source role:} PDF p. 17: deviation from objective priors; PDF p. 19: miscalibrated priors and heterogeneous-prior questions.\par \url{https://bpb-us-w2.wpmucdn.com/voices.uchicago.edu/dist/f/575/files/2021/07/Innsbruck2021_slides.pdf}\par\smallskip
}

\subsection*{Common conventions: Source mathematical conventions}

Unless an entry states otherwise, agent and firm sets are finite. State and action spaces are measurable, strategies and outcome maps are measurable, probability laws are specified, and expectations exist for the targets under discussion. Infinite-horizon discount factors lie in $(0,1)$ and discounted objectives are finite. Norms, tolerances and complexity input sizes must be specified for computational comparisons. Constants or primitives that appear locally are inputs, not empirical facts asserted by this catalogue.

\textbf{Identification.} A statistical model $m\in\mathcal M$ implies an observable law $P_m$ and target $\theta(m)$. At law $P$, its identified set is
\[
\Theta(P;\mathcal M)=\{\theta(m):m\in\mathcal M,\ P_m=P\}.
\]
Point identification holds when this set is a singleton, or equivalently when $P_{m_1}=P_{m_2}$ implies $\theta(m_1)=\theta(m_2)$ for all compatible models. Sharp bounds mean bounds attained, or approached when the boundary is not attained, by compatible models. Writing this set is a definition, not a solution: the substantive task is to establish informative restrictions and derive or estimate the set. An unrestricted-confounding model may give only trivial bounds.

\textbf{Inference.} Identification, estimability and valid uncertainty quantification are separate questions. Fix $\alpha\in(0,1)$. A confidence procedure $C_n$ over a declared law class $\mathcal P$ has asymptotic uniform coverage at level $1-\alpha$ when
\[
\liminf_{n\to\infty}\inf_{P\in\mathcal P}\Pr_P\{\theta(P)\in C_n\}\geq1-\alpha.
\]
For set-identified targets the coverage event must be defined separately, for example coverage of the full identified set. If an entry uses identification-indexed models instead of uniquely indexed laws, the same coverage convention is applied over those models. Pointwise limits do not establish uniform validity, and asymptotic validity does not guarantee finite-sample precision. Sample sizes, dependence structures and weak-identification sequences are part of each application.

\textbf{Counterfactuals.} $Y_i(a)$ is a potential outcome under a specified intervention, including its timing, implementation and versions. It is not $Y_i$ conditional on voluntarily choosing $a$. A full intervention vector permits interference; shorthand scalar treatment notation does not silently impose no interference. Assignment, selection, attrition, anticipation and measurement must be specified. A claim about an instrument's local effect is not automatically a population policy effect.

\textbf{Economic equilibrium.} A counterfactual model includes best responses, constraints, feasibility, market clearing, state transitions and expectations consistent with the stated information. When several equilibria exist, either a declared selector is an input or the target is set-valued. Comparisons across policy regimes must use the stated selection convention. Reduced-form relationships are not presumed invariant after an equilibrium-changing intervention.

\textbf{Optimization and welfare.} Optimization is over the stated feasible set. If attainment is not established, $\sup$ or $\inf$ and $\varepsilon$-optimal choices replace an asserted maximizer or minimizer. For example, with value $V=\sup_{a\in\mathcal A}W(a)<\infty$, an $\varepsilon$-optimal set is $\{a\in\mathcal A:W(a)\geq V-\varepsilon\}$ for $\varepsilon>0$. Benefits, costs, risk and health or environmental outcomes can be aggregated only using declared units and conversion weights. Interpersonal utility weights, the treatment of fiscal transfers and foreign profits, discounting, and the behavioral welfare criterion are explicit inputs. A comparative-static derivative additionally requires existence and appropriate differentiability of the selected counterfactual.

\textbf{Sources.} Local labels [S1], [S2], and so forth restart in every question. Locators identify the relevant abstract, model, section or result at the level actually checked; a bibliographic or abstract-level verification is not represented as inspection of an entire proof. Working papers are distinguished from later publications and changing versions. Source summaries are paraphrased. All URL checks and editorial formulations refer to the audit date stated above.

\subsection*{Common conventions: Additional-source status and mathematical conventions}

\label{audited:conventions}

Of the 100 supplied ECON entries, 65 distinct problems appear here; 32 close
matches refine earlier entries, and three duplicate questions map to existing
entries. Appendix~\ref{app:audited-crosswalk} lists every destination.
The attachment's P/R/S/U distinctions and source audits are retained. Its
precise mathematical formalizations are editorial unless the individual audit
says otherwise. Candidate subsidy targets ECON-004--006 retain their U flags;
the resolved approval-core question ECON-007 updates OP-0202 above.
The following supplied conventions govern both this part and the attributed
ECON refinements elsewhere in the catalogue, subject to local restrictions.

\textbf{Parameterized questions.} A problem family lists inputs that must be fixed before an instance is evaluated: population, horizon, policies, information, data law, model class and welfare weights. These are required choices, not quantities the citations are assumed to specify. Undefined applications should not be treated as identified estimands.

\textbf{Indices and integrability.} Sets of agents, firms and regions are finite unless a population measure is explicitly used. \(i,j\) index units, \(r\) regions, \(t\) dates and \(h\) horizons. \(\mathbb E\) and \(\Pr\) refer to the declared population and law; \(\mathbf1\{A\}\) is an indicator. Every displayed expectation and discounted sum needs existence in the stated domain. Infinite horizons use a discount in \((0,1)\) and a convergence condition; finite horizons may use discount one.

\textbf{Optimization and existence.} A displayed \(\max\), \(\min\), or \(\arg\max\) is conditional on attainment. For finite feasible menus this is immediate; general spaces need continuity and compactness/coercivity or another existence argument. Without attainment, the target is the supremum/infimum and arbitrarily near-optimal feasible choices. \(\inf\varnothing=+\infty\). Empty feasibility and an unidentified target are different issues.

\textbf{Potential outcomes.} \(Y_i(z)\) is the outcome under a specified intervention, not the observed conditional mean among people choosing \(z\). In binary comparisons \(\Delta Y_i=Y_i(1)-Y_i(0)\). Specify assignment, treatment versions, compliance, information and observation horizon. Interference is allowed under a full policy vector; compressed exposure notation needs a justified mapping. No interference, consistency and overlap are substantive assumptions, not automatic consequences of notation.

\textbf{Populations and observation.} Fix baseline eligibility and population weights. Conditioning on post-policy employment, survival, relocation or program uptake can change the population being compared. Death is not ordinary loss to follow-up; outcomes truncated by death need a composite convention, joint survival target, or explicitly defined survivor stratum. Fertility-changing interventions need a population functional rather than an unexplained fixed descendant cohort. Measurement and missingness models must be supplied.

\textbf{Identification.} A structural model \(m\in\mathcal M\) implies observable law \(P_m\) and target \(\theta(m)\). Identification means
\[
 P_{m_1}=P_{m_2}\ \Longrightarrow\ \theta(m_1)=\theta(m_2)
 \quad\text{for all }m_1,m_2\in\mathcal M .
\]
Without identification, the target is
\[
 \Theta(P;\mathcal M)=\{\theta(m):m\in\mathcal M,\ P_m=P\}.
\]
Writing \(\theta(P)\) before establishing identification can hide incompatible latent models. Confidence sets additionally need a sampling law, confidence level, asymptotic sequence and uniformity class. Point identification, coverage and informative precision are separate properties.

\textbf{Mediation.} Total effects of randomized assignment are distinct from cross-world natural indirect effects. Belief/effort-channel effects require a meaningful mediator intervention and additional measurement, exclusion or mediation assumptions. Randomization of the initial policy alone does not supply those assumptions. Report bounds or interventional alternatives when the stronger target is unsupported.

\textbf{Equilibrium.} A counterfactual \(W(z)\), \(p(z)\), or \(\mu_t^z\) requires individual optimization, feasibility, government/resource budgets, market clearing and state-transition laws. State equilibrium existence and selection, or report ranges over compatible equilibria. Initial-state integration includes subsequent endogenous transitions; current-state integration must use the policy-dependent distribution. Reduced-form invariance after policy is not assumed.

\textbf{Welfare accounts.} Scores, utility, health, dollars and ecological quantities require explicit conversion before addition. Fees, taxes, subsidies, wages and extortion payments are transfers between parties; distributional weights can make them welfare relevant, but their face value is not automatically a resource loss. State ownership, geographic boundaries, foreign-income weights and fiscal finance. Do not add profits to household welfare already including dividends, or count land capitalization and the underlying benefit twice. A benefit-cost expression must distinguish gross benefits from net welfare and subtract every real cost exactly once.

\textbf{Derivatives and risk.} Log derivatives require positive arguments; local responses require admissible perturbations and specified equilibrium branches. Differentiation under expectations needs regularity/domination, and boundaries may allow only one-sided derivatives. Risk uses a specified law; ambiguity uses a declared nonempty law set on a common history space. Adaptive policies must be nonanticipative. A robust criterion is an editorial choice unless attributed explicitly.

\textbf{Scope overlaps.} Allocation, collective choice, contracts, household bargaining and coordinated policy overlap economics with optimization and game theory. They are retained because they appeared in the original catalogue; no claim of a strictly game-theory-free selection is made.
% END ECONOMICS STATEMENT
\end{document}
